\section{Geometry Generates Gauge}

Geometry enters this chapter as boundary grammar, not as a container waiting
behind the record. A geometry is a stabilized way of saying which separations
matter, which boundaries have been committed, which interiors remain
unread, and which transports preserve distinguishability. When such a
boundary grammar is asked to carry the charged pair residue, it cannot remain
merely spatial. It must generate a gauge: a rule for reading how the
distinction is transported, compared, and discriminated.

The generating step begins with a boundary depth. An interior reading may
carry no charge when read alone. A boundary reading may carry a unit
obligation because it is where the interior must answer to an exterior
comparison. The grammar distinguishes these two faces. It forbids the
interior from inventing charge without boundary support, and it forbids the
boundary from acting as a decorative shell. The boundary is the place where
unspent orientation becomes readable.

A compact-disc threshold gives a modest illustration. The read head does not
begin with a metric as a metaphysical object. It admits a new mark only when
two surface configurations are separated enough to force a new event in the
record. The gauge of separation is inferred from admissible distinction. In
the chapter's grammar, spatial boundary works similarly. A boundary is not a
line drawn around a pre-given space. It is the rule that says when a
separation has become expensive enough to enter the ledger.

Once the boundary carries this rule, it can be compiled into an action. The
word compiled should be read grammatically: boundary commitments are gathered
into a finite reader that compares interior and exterior charges, preserves
the anchor, and tests the coupled residue. The resulting action is not added
beside geometry. It is geometry's boundary obligation expressed as a gauge
test. That test is precisely what the electron section called the
discriminating action.

This is the sense in which geometry generates gauge. The grammar starts with
boundary separations and asks what must be carried so that a pair residue can
be read without contradiction. It must carry a baseline. It must carry
orientation. It must carry the distinction between separated parts and the
coupled residue. It must carry the boundary debt that the interior alone
cannot settle. Those carriers are gauge carriers. The gauge is the transport
grammar forced by the boundary.

The Aharonov--Bohm illustration again shows the shape. A region may have no
local field reading along two branches, yet the closed comparison can return
a phase difference. The physically relevant distinction is then not the
local force reading alone but the connection that carries memory around the
loop. In this chapter the point is more basic: a boundary grammar that
preserves distinguishability around a comparison must become a gauge grammar
when the transported sign matters.

This conversion also protects the electron reading from arbitrariness. If
the discriminating action were merely chosen, the \(-1\) reading would look
like a convention. But when boundary geometry generates the action, the
charge reading is forced by the boundary's own obligation to reconcile
interior and exterior. The gauge reads what the geometry has made
unavoidable. It discriminates because the boundary has produced a residue
that cannot be quotiented away.

The white-hole image can be used carefully here. A source boundary exports a
reading that the interior cannot absorb as its own private fact. The export
does not mean a mythic object has been placed at the boundary. It means the
boundary grammar carries an asymmetry: the interior may be charge-silent,
while the boundary carries the unit condition needed for discrimination.
Gauge is the rule that transports that boundary condition into a readable
action.

Thus geometry generating gauge is not a slogan of unification. It is a
discipline of origin. The grammar permits geometry as stabilized boundary
bookkeeping. It forbids geometry from claiming charged transport until a
residue demands it. It then identifies the necessary transport rule as gauge.
When the gauge acts on the pair residue, the electron reading is not an
extra particle inserted into geometry. It is the boundary-generated action's
negative charge.

The section also prepares the split. If gauge is generated by geometry, it
is tempting to add the geometric and gauge readings together as if they were
two independent components of one finite vector. The grammar will forbid
that. Generation is not direct summation. The gauge is born from the boundary
obligation of geometry, but once the reading is made, the finite carrier must
respect the difference between the geometric face and the gauge face.

